\section{The same depth obstruction for ReLU}
\label{sec:relu-depth-extension}
The positive-label memorizer of Section~\ref{sec:relu-near-baseline} yields a ReLU version with a fixed width-two target. Shifting the hinge is useful: the unshifted function $\operatorname{ReLU}(T)$ is affine on this augmented four-point support, whereas the shifted hinge below has a nonzero residual from all affine input functions.

\begin{theorem}[ReLU depth lower bound]
\label{thm:relu-depth-lower}
There is an absolute $\eta>0$ such that, for every $\kappa>0$, there are finite $C_\kappa,H_\kappa$ with the following property. For every integer $H\ge H_\kappa$, put
\[
 D=\lfloor\eta H\rfloor,\qquad d=D+2,\qquad n=4\cdot2^D,
\]
the uniform input law on $\mathcal X_D$ with
\[
 g(x)=\frac12+\frac14\operatorname{ReLU}(T-\tfrac12),\qquad
 Y=g(X)+\epsilon/8,\quad \Pr(\epsilon=\pm1)=1/2,
\]
where the signs are independent of the inputs, has the following property. With probability at least $9/10$ over $n$ IID observations, simultaneously for every
\[
 N\ge e^{C_\kappa H},\qquad 0<s\le N^{-1-\kappa},\qquad \lambda>0,
\]
the ordinary raw posterior mean of the ReLU learner has risk at least
$1/64+2^{-22}$. The same conclusion holds for clipping its mean, the mean of clipped draws, arbitrary data-dependent mixtures of these predictions, and clipping after mixing.

If the learner is additionally given any independent centered Gaussian affine output component $b^\top(1,x)$ with finite positive semidefinite covariance, the conclusion remains valid for ordinary raw posterior means and arbitrary mixtures of them, with excess lower bound $2^{-40}$ in place of $2^{-22}$. This statement is simultaneous over all such affine covariances.

The target has width two and coefficient bound one at every hidden depth, independently of $d$ and $H$.
\end{theorem}
\begin{proof}
The target takes values $(1/2,1/2,1/2,9/16)$ on the ordered $T$ values. Thus all labels lie in $[3/8,11/16]$, and all empirical atom means are positive. Assign any value in $[0,1]$ at unobserved atoms. The ReLU near-baseline theorem supplies an exact memorizer with energy $K\le(1+\delta)H+c$ and an unchanged-covariance Gaussian comparison whose RMS error is arbitrarily small uniformly over the stated widths and scales. The resulting comparison loss is $r_0+\zeta$: within-atom residuals have zero sum, so their cross term vanishes.

ReLU is zero-origin and one-Lipschitz, so the same operator-energy floor and entropy inequality apply. Its native prior output variance on this support is exactly $\frac2d s(Ns/2)^H\le s$, giving the same zero comparison. Choose the fixed norm threshold, entropy slack, memorizer energy slack and comparison error sufficiently small. For $0<\lambda\le16/H$ the objective correction is
\[
 s+\lambda s\mathcal B_a=O_a(s+Ns),
\]
which tends uniformly to zero at the stated widths. The segment lemma therefore puts each raw mean and each mean of clipped draws within any prescribed fixed tolerance of $\{pb:0\le p\le1\}$. For $\lambda\ge16/H$, the posterior operator-energy tail proof applies without change: it uses only the zero comparison, the Gaussian operator-norm moment and the zero-origin Lipschitz output-product bound. Bounded hidden coordinates are unnecessary. Since $Ns\to0$ uniformly, it makes all those mean conventions pointwise exponentially small in $H$. Clipping a mean preserves its distance from the segment.

On the occupancy event, $\|g(S)\|_n^2\ge1/4$ and singleton noise gives
$\|b-g(S)\|_n^2\ge1/192$. Also $D_S\le1$ and
$|b-2g(S)|\le1$ pointwise. Conditional Hoeffding gives
$|\langle b-g(S),g(S)\rangle_n|\le1/2048$
except on an exponentially small failure event. Choose the segment tolerance at most $1/2048$. The common witness is then at least
\[
 1/256-1/2048-1/2048\ge1/512.
\]
The population witness $h_S(x)=\sum_{i:X_i=x}(Y_i-2g(x))$ has squared norm at most four when $\sum_xk_x^2/n\le4$. Cauchy--Schwarz gives excess at least $2^{-20}$, leaving the stated slack. The same segment and witness survive mixtures and clipping after mixing.

For the affine extension, let $\Phi$ have rows $(1,X_i)^\top$, let $\Sigma$ be the affine covariance and set
\[
 B=(I+\Phi\Sigma\Phi^\top/(n\lambda s))^{-1}.
\]
If $P$ projects orthogonally away from the full empirical affine span, then
$0\prec B\preceq I$ and $PB=BP=P$. Gaussian integration gives neural marginal loss $\|B^{1/2}(y-z)\|_n^2$, and the total mean $f$ satisfies
\[
 f=(I-B)y+Bm_z,\qquad Pf=Pm_z.
\]
Within-atom residuals are affine-orthogonal. Apply the segment lemma in the $B$ inner product; its floor follows from $\|B^{1/2}z\|\le\|z\|$. Both prior comparisons and the native memorizer remain valid because $B\preceq I$. Thus
$\|P m_z-pP b\|_n$ is arbitrarily small uniformly in $\Sigma$ at low temperatures. At high temperatures the operator-energy tail still gives pointwise smallness of the neural mean: its zero comparison is bounded by $\widehat R_B(0)+s\le1+s$. No bound on $\Sigma$ is used.

The population affine residual of $g$ on the four $T$ values is
\[
 \frac1{320}(-1,3,-3,1),\qquad
 \|g-\operatorname{proj}_{\rm affine}g\|_2^2=\frac1{20480}.
\]
Indeed the two even directions are spanned by $1$ and $\sqrt{3/4-T^2}$, and the only odd affine direction is $T$. The displayed residual is odd and orthogonal to $T$. Its corresponding affine fitted values are
$(161,157,163,179)/320$, all below $3/4$.
Use the whitened affine coordinates $(1,T,z,\sqrt{3/4-T^2})$ as in the tanh construction. They span all affine input functions, and their Gram matrix has a fixed positive lower eigenvalue. Bounded-entry concentration at precision $O((D+3)^{-2})$ therefore gives
\[
 \|Pg(S)\|_n^2>2^{-16}
\]
with failure probability tending to zero. Conditional on inputs, the affine projection of the independent sign noise has expected squared norm at most $(D+3)/(64n)$. Markov and the singleton event give $\|P(b-g)\|_n^2\ge1/384>2^{-16}$ with probability tending to one. Conditional Hoeffding makes the projected alignment at most $2^{-20}$ in absolute value. The projected witness has norm at most one.

Make the projected segment tolerance at most $2^{-20}$.
The common empirical separator is then at least $2^{-17}$.
Let $\widehat\beta$ be the empirical affine coefficient fitting $y-2g(S)$.
Its population fitted values are the negatives of the bounded affine values above. The same bounded Gram and noise-moment concentration gives
$\sup_{x\in\mathcal X_D}|(1,x)^\top\widehat\beta|\le1$.
For completeness, entry errors $O((D+3)^{-2})$ give coefficient error $O((D+3)^{-1})$; multiplying by the support norm $O(\sqrt{D+3})$ leaves a vanishing uniform error. The failure probability is at most a polynomial in $D$ times $e^{-cn/(D+3)^4}$.

The population witness
\[
 h(x)=\sum_{i:X_i=x}\{Y_i-2g(x)-(1,x)^\top\widehat\beta\}
\]
pairs with every $f-g$ exactly as the empirical projected witness.
Each summand has absolute value at most two; the count event therefore gives $\|h\|_2^2\le16$. The excess is at least $2^{-38}$, implying the stated $2^{-40}$. This witness is common to all temperatures, widths, scales and affine covariances, hence to their arbitrary mixtures. Clipping an affine prediction is not included in this projected argument.

All event failures tend to zero because $n$ is exponential while the affine dimension is linear in $H$. Increase the depth threshold to make their sum at most $1/10$. Finally, the target is implemented by first-layer constant one and $\operatorname{ReLU}(T-1/2)$, followed by identity relays and readout $(1/2,1/4)$; all coefficients have magnitude at most one.
\end{proof}
